\chapter{Intervention Design Patterns}
\label{ch:interventions}

\section{From Description to Intervention}

The previous chapters argued that behavioural failure is often a state-transition
problem rather than a simple shortage of desire. That claim matters only if it
changes design practice. An intervention chapter therefore has a stricter job
than a motivational chapter: it must specify which control variable is being
changed, why that variable is cheaper or more reliable than direct exhortation,
and which parts of the design are supported by stronger evidence rather than
mere intuition.

In this book, intervention design is governed by five ideas:
\begin{enumerate}[nosep]
  \item interventions act on transitions, not on abstract identity claims;
  \item decision windows matter because leverage is timing-sensitive;
  \item repeated exposure and history dependence matter because frameworks are
        not rebuilt from zero at every attempt;
  \item attractor and metastability language helps explain why some local
        redesigns have outsized downstream effects;
  \item source discipline matters because empirical findings, theoretical
        mechanism-building, and speculative frontier claims should not be mixed
        as if they had equal status.
\end{enumerate}

\section{Evidence Lanes For Intervention Thinking}

This chapter uses the book's evidence policy explicitly.

\begin{itemize}[nosep]
  \item \textbf{Strong lane.} State-transition evidence from consciousness and
        neural-dynamics work supports the general claim that systems can occupy
        distinct global regimes and move between them through structured
        transitions. This supports treating intervention as regime-shaping
        rather than mere self-talk.
  \item \textbf{Medium lane.} Metastability, attractor, integrator, causal
        emergence, and nonlinear-control perspectives support the design idea
        that repeated local inputs can stabilize or redirect future behaviour.
  \item \textbf{Speculative lane.} Frontier consciousness-physics proposals may
        inspire hypotheses about large-scale organization, but they do not
        justify practical intervention rules in this chapter.
  \item \textbf{Design inference.} The concrete protocols below are applied
        synthesis. They are informed by the strong and medium lanes, but they
        remain design constructions rather than direct experimental theorems.
\end{itemize}

\begin{proposition}[Interventions should target the transition geometry]
When a failure is recurrent, the first design question should usually be:
which alteration most cheaply changes the probability of entering the target
state during a real decision window?
\end{proposition}

\begin{discussion}
This proposition is stronger than ``make the task easier.'' It says the
intervention should be evaluated in transition terms. If the current state is a
locally stable basin, then a useful intervention either lowers the barrier,
changes the timing of attempted movement, or alters the local environment so
that the old state becomes less self-restoring. That logic is consistent with
the book's state-transition framing and with the broader dynamical-systems
language of attractors, metastability, and bifurcation-like regime shifts.
\end{discussion}

\section{Comparative Intervention Map}

\begin{itemize}[nosep]
  \item \textbf{Window capture.} Target: timing and entry leverage.
        Evidence lane: strong + design inference.
        Use: move the attempt to a boundary such as wake-up, post-break
        re-entry, or commute transition.
  \item \textbf{State preloading.} Target: switching cost and path clarity.
        Evidence lane: medium + design inference.
        Use: prepare the exact next page, sentence, route, or tool before the
        window opens.
  \item \textbf{Environmental re-weighting.} Target: cue landscape and local
        reinforcement.
        Evidence lane: strong/medium + design inference.
        Use: remove cues that rebuild the old state and add cues that make the
        target state easier to enter.
  \item \textbf{History shaping.} Target: emergence and path dependence.
        Evidence lane: medium + design inference.
        Use: repeat the same entry move in the same context so local history
        accumulates.
  \item \textbf{Minimum viable entry.} Target: barrier crossing under limited
        will.
        Evidence lane: medium + design inference.
        Use: define the smallest real action that still counts as entering the
        target regime.
  \item \textbf{External scaffolding.} Target: effective control input.
        Evidence lane: medium + design inference.
        Use: use schedules, partners, or deadlines when private control is too
        noisy to trigger reliable entry.
\end{itemize}

\section{Design Principles}

\begin{proposition}[Cheapest controllable variable principle]
When several barrier terms are elevated, the first intervention should target
the cheapest controllable variable that materially changes transition odds,
rather than the most morally impressive variable.
\end{proposition}

\begin{discussion}
This principle follows the logic of control design. In a recurrent failure,
raising internal resolve may be possible, but it is usually expensive, noisy,
and hard to verify. By contrast, moving the phone away, pre-marking the next
problem, or attaching the task to an already recurring boundary changes the
state-transition conditions directly and cheaply. In dynamical-systems terms,
small local manipulations can matter if they change the system near a sensitive
entry point.
\end{discussion}

\begin{proposition}[History-dependent reinforcement principle]
An intervention should be judged not only by whether it succeeds once, but by
whether it makes the next attempt easier through accumulated local history.
\end{proposition}

\begin{discussion}
This principle is where emergence enters practice. If a reader repeatedly uses
the same location, same startup cue, and same minimum viable entry action, then
future transitions inherit a partially trained pathway. That is not yet a proof
of full attractor formation in the strict mathematical sense, but it is a
reasonable design inference from medium-lane ideas about metastability, causal
emergence, and path dependence: repeated local structure can produce more
stable higher-level patterns.
\end{discussion}

\section{Implementation Template}

\begin{templatebox}
\textbf{Research-backed intervention sheet}
\begin{enumerate}[nosep]
  \item Current state:
  \item Target state:
  \item Dominant barrier term ($\Reward$ / $\Switch$ / $\Env$ / timing / ambiguity):
  \item Decision window to exploit:
  \item Which evidence lane most directly supports this intervention choice?
        (`strong` / `medium` / `design inference`)
  \item One environmental re-weighting move:
  \item One state-preloading move:
  \item One minimum viable entry action:
  \item One history-shaping rule to repeat across sessions:
  \item Verification rule after the session:
\end{enumerate}
\end{templatebox}

\section{Worked Implementation}

\begin{example}[Rebuilding post-dinner study entry as a transition design]
Consider a reader who repeatedly intends to do technical reading after dinner
but instead enters a phone-browsing state that lasts the rest of the evening.

\textbf{Step 1: Diagnose the failure.}
\begin{itemize}[nosep]
  \item Current state: post-dinner browsing with rapid novelty and low effort.
  \item Target state: twenty minutes of annotated reading with one written
        question.
  \item Dominant barrier terms: high current reward, high ambiguity about the
        first step, and a narrow post-cleanup window.
\end{itemize}

\textbf{Step 2: Identify the evidence lane.}
\begin{itemize}[nosep]
  \item Strong lane support: state-transition research justifies treating the
        dinner-to-evening boundary as a meaningful regime-change opportunity.
  \item Medium lane support: metastability and attractor-style reasoning
        support the expectation that repeated use of the same cue and same desk
        setup can make future entry more reliable.
  \item Design inference: the exact configuration below is an applied design,
        not a direct citation result.
\end{itemize}

\textbf{Step 3: Redesign the transition.}
\begin{enumerate}[nosep]
  \item \textbf{Window capture:} define the intervention point as the two
        minutes immediately after dish cleanup.
  \item \textbf{Environmental re-weighting:} phone leaves the room before the
        meal starts, so the old evening attractor is harder to reconstruct.
  \item \textbf{State preloading:} the book is opened in advance with a sticky
        note stating one question to answer.
  \item \textbf{Minimum viable entry:} read one page and write one margin note.
  \item \textbf{History shaping:} use the same chair, same lamp, and same first
        notebook action for seven consecutive sessions.
  \item \textbf{Verification:} record whether the first margin note appears
        before 8:15 PM.
\end{enumerate}

\textbf{Why this is better than generic advice.}
The design does not ask the reader to feel more serious. It changes the local
geometry of the transition. The dinner boundary becomes a designed decision
window; the target state is partially built before entry; the old state is made
less recoverable; and repeated history is used to increase the odds that the
next evening begins from a more favourable local configuration.
\end{example}

\section{Failure Review Protocol}

When an intervention fails, the review should remain diagnostic and source-aware:
\begin{enumerate}[nosep]
  \item Was the chosen decision window real, or was it only imagined after the
        fact?
  \item Did the intervention change a controllable variable, or did it merely
        restate the goal?
  \item Was the target state sufficiently preloaded to reduce switching cost?
  \item Did the environment still recreate the old local attractor too quickly?
  \item Was the minimum viable entry still too large for the available
        volitional budget?
  \item Did the design include any history-shaping repetition, or was each
        attempt effectively a fresh start?
  \item Is the intervention justified by a strong or medium lane, or is it only
        a speculative guess that should be downgraded?
\end{enumerate}

\section{Recommended Reading}

\begin{readingbox}
\textbf{Foundation}
\begin{itemize}[nosep]
  \item Eugene Izhikevich, \emph{Dynamical Systems in Neuroscience}, for a
        rigorous but usable route from nonlinear state dynamics to intervention
        thinking.
  \item Steven Strogatz, \emph{Nonlinear Dynamics and Chaos}, for intuition
        about regime change, local stability, and why small changes near a
        boundary can matter.
  \item John Holland, \emph{Emergence}, for the idea that repeated local rules
        can accumulate into durable higher-level structure.
\end{itemize}
\textbf{Modern Research}
\begin{itemize}[nosep]
  \item \emph{Dynamical structure-function correlations provide robust and
        generalizable signatures of consciousness in humans} (Communications
        Biology, 2024), as support for treating conscious states as structured
        dynamical regimes rather than undifferentiated ``motivation.''
  \item \emph{Cortical connectivity, local dynamics and stability correlates of
        global conscious states} (Communications Biology, 2025), for the link
        between stability structure and global state organization.
  \item \emph{Falling asleep follows a predictable bifurcation dynamic}
        (Nature Neuroscience, 2025), for a concrete example of state transition
        analysis that strengthens window-sensitive intervention thinking.
  \item \emph{Metastability demystified: the foundational past, the pragmatic
        present and the promising future} (Nature Reviews Neuroscience, 2024),
        for the medium-lane vocabulary behind temporary stabilization and
        flexible switching.
  \item \emph{Associative conditioning in gene regulatory network models
        increases integrative causal emergence} (Communications Biology, 2025),
        as a bridge source for history dependence and emergence-informed
        intervention design.
\end{itemize}
\textbf{Frontier}
\begin{itemize}[nosep]
  \item Frontier consciousness-physics proposals should be read only as
        hypothesis-generating extensions. They may inspire alternative system
        pictures, but they do not currently justify operational intervention
        rules in this chapter.
\end{itemize}
\end{readingbox}

\begin{summarybox}
This chapter turns the book's research spine into operational design. The main
practical lesson is not ``try harder.'' It is: diagnose the transition, choose
the cheapest controllable variable, exploit a real decision window, and prefer
designs that create better future history rather than one-off acts of effort.
The strong lane justifies transition-sensitive thinking; the medium lane helps
explain why repetition and local structure matter; the final protocol remains a
design inference whose value must be checked in use.
\end{summarybox}
